\documentclass[11pt,a4paper]{article}

\usepackage[utf8]{inputenc}
\usepackage[T1]{fontenc}
\usepackage{amsmath, amsthm, amssymb, mathtools}
\usepackage{microtype}
\usepackage[margin=1in]{geometry}
\usepackage[dvipsnames]{xcolor}
\usepackage{hyperref}
\usepackage{booktabs}
\usepackage{enumitem}

\hypersetup{
  colorlinks=true,
  linkcolor=blue!50!black,
  citecolor=blue!50!black,
  urlcolor=blue!50!black,
}

\newtheorem{theorem}{Theorem}
\newtheorem{lemma}[theorem]{Lemma}
\newtheorem{corollary}[theorem]{Corollary}
\theoremstyle{definition}
\newtheorem{remark}[theorem]{Remark}
\newtheorem{definition}[theorem]{Definition}

\newcommand{\C}{\mathbb{C}}
\newcommand{\R}{\mathbb{R}}
\newcommand{\Z}{\mathbb{Z}}
\newcommand{\conj}[1]{\overline{#1}}
\DeclareMathOperator{\re}{Re}
\DeclareMathOperator{\im}{Im}
\DeclareMathOperator{\Arg}{Arg}
\DeclareMathOperator{\Log}{Log}

\title{Reflection-symmetric discs:\\
a short certificate that the zeros of $\zeta$ below height $2000$\\
are simple and lie on the critical line}
\author{Jude Gomila}
\date{Draft of 20 November 2026}

\begin{document}
\maketitle

\begin{abstract}
A disc centred on the critical line is mapped to itself by the reflection
$s \mapsto 1-\conj{s}$. If such a disc contains exactly one zero of the
Riemann zeta function, counted with multiplicity, that zero is simple and
lies on the critical line. We combine this elementary observation with
winding numbers computed in ball arithmetic to certify that $\zeta$ has
exactly $3034$ zeros in the rectangle $[-1,2]\times[-2000,2000]$, all simple
and all on the line $\re s=\tfrac12$. The argument uses no Gram points, no
Turing bound and no properties of Hardy's function: its correctness reduces to
rigorous enclosures of $\zeta$ on small boxes, however those are computed. The verified
height is far below what is known; the contribution is a short, self-contained
and quickly reproducible argument, together with one independent certificate
per zero. Nothing here bears on zeros above height $2000$.
\end{abstract}

\section{Introduction}

Let $\zeta$ be the Riemann zeta function and write $s=\sigma+it$. The Riemann
hypothesis asserts that every zero of $\zeta$ in the strip $0<\sigma<1$ has
$\sigma=\tfrac12$. It has been verified numerically to great height: Platt and
Trudgian \cite{PlattTrudgian} prove that all zeros with $0<t\le 3\cdot 10^{12}$
lie on the critical line. Such verifications follow a method
going back to Turing \cite{Turing} and developed by Brent \cite{Brent} and
others: sign changes of Hardy's function $Z(t)$ locate zeros on the line, and
a bound on the argument of $\zeta$ shows that no others exist.

This note records a different and much more modest verification. It reaches
only height $2000$, and it is not competitive in cost. What it offers is
brevity. The whole argument is the argument principle applied to two kinds of
contour, discs and one rectangle, with every winding number obtained from
rigorous enclosures of $\zeta$ itself.

\begin{theorem}\label{thm:main}
Let $R=[-1,2]\times[-2000,2000]$. The function $\zeta$ has exactly $3034$
zeros in $R$, counted with multiplicity. Each of them is simple and has real
part $\tfrac12$. In particular $\zeta$ has exactly $1517$ zeros with
$0<\im s<2000$, and all of them are simple and lie on the critical line.
\end{theorem}

The proof is computer-assisted. Section~\ref{sec:lemma} proves the lemma that
turns a count into a location. Section~\ref{sec:winding} explains how a
winding number is certified. Section~\ref{sec:cert} states the three computed
facts and derives Theorem~\ref{thm:main} from them. Section~\ref{sec:trust}
lists exactly what must be trusted, and Section~\ref{sec:limits} states the
limits of the method.

\paragraph{What is and is not new.}
The statement of Theorem~\ref{thm:main} is a very special case of
\cite{PlattTrudgian}. The lemma of Section~\ref{sec:lemma} is elementary and
we make no claim of priority for it. The point of the note is the combination:
a verification whose logic fits on a page, whose data are $1517$ discs and one
rectangle, and which runs in about three minutes. The discs also serve a
second purpose. Each is a certificate for one zero that can be checked on its
own, and the interactive viewer that accompanies this note displays a zero as
certified only when its disc matches a recorded one exactly.

\section{The trap lemma}\label{sec:lemma}

Write $\rho(s)=1-\conj{s}$ for the reflection in the line $\re s=\tfrac12$.
Its fixed points are exactly the points of that line.

\begin{lemma}[Trap lemma]\label{lem:trap}
Let $D$ be a closed disc centred at a point of the line $\re s=\tfrac12$.
Let $F$ be holomorphic on a neighbourhood of $D$ and satisfy
\begin{equation}\label{eq:sym}
  F(1-\conj{s})=\conj{F(s)}
\end{equation}
there. If $F$ has exactly one zero in $D$, counted with multiplicity, then
that zero is simple and lies on the line $\re s=\tfrac12$.
\end{lemma}

\begin{proof}
The reflection $\rho$ maps $D$ onto itself. By \eqref{eq:sym}, if $F(s_0)=0$
then $F(\rho(s_0))=0$. Moreover the two zeros have the same multiplicity:
\eqref{eq:sym} says $F=c\circ F\circ\rho$ with $c$ complex conjugation, and
composition with the anti-holomorphic bijections $\rho$ and $c$ preserves the
order of vanishing. If $s_0\in D$ were a zero with $\re s_0\neq\tfrac12$, then
$\rho(s_0)\neq s_0$ would be a second zero in $D$, and $D$ would contain at
least two zeros counted with multiplicity. Hence the single zero is fixed by
$\rho$, so it lies on the line. A count of one means multiplicity one.
\end{proof}

The lemma applies to the completed function
\[
  \xi(s)=\tfrac12\,s(s-1)\,\pi^{-s/2}\,\Gamma\!\left(\tfrac{s}{2}\right)\zeta(s),
\]
which is entire, satisfies $\xi(s)=\xi(1-s)$, and is real on the real axis, so
that $\xi(\conj{s})=\conj{\xi(s)}$. Together these give \eqref{eq:sym} for
$F=\xi$.

\begin{corollary}\label{cor:zeta}
Let $D$ be a closed disc centred at $\tfrac12+it_0$ with $t_0>0$ and radius
$r<t_0$. If $\zeta$ has exactly one zero in $D$, counted with multiplicity,
then that zero is simple and lies on the critical line.
\end{corollary}

\begin{proof}
Every point of $D$ has positive imaginary part. The factor
$\tfrac12 s(s-1)\pi^{-s/2}\Gamma(s/2)$ is holomorphic and non-vanishing there:
its only zeros and poles are at real points. So $\zeta$ and $\xi$ have the
same zeros in $D$, with the same multiplicities, and Lemma~\ref{lem:trap}
applies to $\xi$.
\end{proof}

The lemma converts a statement about \emph{how many} zeros a region holds
into a statement about \emph{where} they are. Counting is what the argument
principle does.

\section{Certified winding numbers}\label{sec:winding}

Let $\Gamma$ be a closed, piecewise smooth, positively oriented simple contour
and $f$ a function meromorphic on a neighbourhood of $\Gamma$ and its
interior, with no zero or pole on $\Gamma$. The winding number
\[
  W(f,\Gamma)=\frac{1}{2\pi}\,\Delta_\Gamma \arg f
\]
equals the number of zeros of $f$ inside $\Gamma$ minus the number of poles,
both counted with multiplicity. We need a way to obtain $W$ exactly from
finitely many rigorous evaluations.

\begin{lemma}[Half-plane lemma]\label{lem:half}
Let $P$ be a path from $a$ to $b$ and $f$ a function continuous on $P$ with
$f(a)\neq0$. Suppose
\begin{equation}\label{eq:half}
  \re\frac{f(z)}{f(a)}>0\qquad\text{for all } z\in P .
\end{equation}
Then $f$ has no zero on $P$, and the continuous change of $\arg f$ along $P$
is
\[
  \Delta_P\arg f=\Arg\frac{f(b)}{f(a)}\in\left(-\tfrac{\pi}{2},\tfrac{\pi}{2}\right),
\]
where $\Arg$ is the principal argument.
\end{lemma}

\begin{proof}
Put $g=f/f(a)$. By \eqref{eq:half}, $g$ maps $P$ into the open right
half-plane, so $g$ and hence $f$ do not vanish on $P$. The principal logarithm
$\Log$ is continuous on that half-plane, so $z\mapsto\im\Log g(z)$ is a
continuous choice of $\arg g$ along $P$. Its total change is
$\Arg g(b)-\Arg g(a)=\Arg g(b)$, because $g(a)=1$. Finally $\arg f$ and
$\arg g$ differ by a constant.
\end{proof}

\paragraph{The procedure.}
Divide $\Gamma$ into consecutive pieces $P_1,\dots,P_n$, with $P_j$ running
from $a_j$ to $b_j=a_{j+1}$ and $b_n=a_1$. If \eqref{eq:half} holds on every
piece, then $f\neq0$ on $\Gamma$ and
\begin{equation}\label{eq:sum}
  W(f,\Gamma)=\frac{1}{2\pi}\sum_{j=1}^{n}\Arg\frac{f(b_j)}{f(a_j)} .
\end{equation}
Both the hypothesis and the sum are evaluated in ball arithmetic. A ball is a
midpoint with a radius, and a ball-valued routine returns a ball guaranteed to
contain the exact value for every point of its input ball.
\begin{enumerate}[leftmargin=2em]
\item For each endpoint, compute a ball $A_j\ni f(a_j)$ and check
  $0\notin A_j$.
\item For each piece, choose a box $B_j\supseteq P_j$ and compute a ball
  $I_j\supseteq f(B_j)$. Check that the ball $I_j/A_j$ lies in the open right
  half-plane. This proves \eqref{eq:half} on $P_j$. If the check fails, bisect
  $P_j$ and repeat on the halves.
\item Sum balls containing $\Arg\!\big(f(b_j)/f(a_j)\big)$ and divide by
  $2\pi$. The result is a ball $T$ that contains $W(f,\Gamma)$.
\item Since $W$ is an integer, if $T$ has radius less than $\tfrac14$ and
  contains the integer $n_0$, then $W=n_0$.
\end{enumerate}
The procedure can fail to terminate at a given precision and bisection depth,
in which case nothing is claimed. It cannot return a wrong integer, provided
the enclosures are correct. In particular the result does not depend on any
sampling density: the half-plane check is made on whole boxes, not on sample
points.

\section{The certificate}\label{sec:cert}

\subsection{The data}

The data are $N=1517$ pairs $(t_k,r_k)$, $t_1<\dots<t_N$, each number an
IEEE double and therefore an exact dyadic rational, and the rectangle
$R=[-1,2]\times[-2000,2000]$. Put
\[
  D_k=\{\,s:\ |s-(\tfrac12+it_k)|\le r_k\,\},\qquad k=1,\dots,N .
\]
The $t_k$ are approximations to the heights of the zeros; the radii are
$r_k=\min\big(0.6,\ 0.4\,g_k\big)$, where $g_k$ is the distance from $t_k$ to
the nearest other $t_j$, except that $r_1=0.35$. The first $25$ pairs are the
ones the accompanying viewer draws. How the centres were found plays no part
in the proof. They are proposals; the proof for each disc is its winding
number.

\subsection{Three computed facts}

\begin{enumerate}[label=\textbf{(C\arabic*)},leftmargin=3em]
\item\label{C1} \textbf{Discs.} For every $k$, $\zeta$ has no zero on
  $\partial D_k$ and $W(\zeta,\partial D_k)=1$.
\item\label{C2} \textbf{Rectangle.} $\zeta$ has no zero on $\partial R$ and
  $W(\zeta,\partial R)=3033$.
\item\label{C3} \textbf{Geometry.} In exact rational arithmetic,
  $t_{k+1}-t_k>r_k+r_{k+1}$ for $1\le k<N$, $\;t_1-r_1>0$,
  $\;t_N+r_N<2000$, and $r_k\le 0.6$ for all $k$.
\end{enumerate}

Facts \ref{C1} and \ref{C2} are obtained by the procedure of
Section~\ref{sec:winding} with $f=\zeta$. For a disc, the circle is cut into
$16$ equal arcs to begin with. An arc of angular fraction $\delta$ on a circle
of radius $r$ has length $2\pi r\delta$, so every point of it is within
$\pi r\delta$ of its midpoint; the box $B_j$ is the square of half-side
$3.1416\,r\delta$ about the midpoint. For the rectangle, the sides are cut into
$8$ (horizontal) or $64$ (vertical) segments to begin with, and each piece is
its own box. Table~\ref{tab:stats} gives the sizes of the computation. The
quantities in \ref{C3} are, to four figures,
\[
  \min_k\big(t_{k+1}-t_k-r_k-r_{k+1}\big)=0.0195,\qquad
  t_1-r_1=13.78,\qquad t_N+r_N=1999.90 .
\]

\begin{table}[ht]
\centering
\begin{tabular}{lrr}
\toprule
 & discs (all $1517$) & rectangle \\
\midrule
winding number & $1$ each & $3033$ \\
boundary pieces after bisection & $517{,}522$ & $65{,}157$ \\
pieces per contour, min / median / max & $82$ / $297$ / $2365$ & --- \\
\bottomrule
\end{tabular}
\caption{Size of the certified computation. Working precision is $128$ bits
and the bisection depth limit is $40$. The whole run takes about $200$ seconds
on one core of a laptop.}
\label{tab:stats}
\end{table}

\subsection{Proof of Theorem~\ref{thm:main}}

\emph{One zero per disc.} The only pole of $\zeta$ is at $s=1$, which is not in
any $D_k$ because every point of $D_k$ has imaginary part at least
$t_1-r_1>0$ by \ref{C3}. So $\zeta$ is holomorphic on a neighbourhood of
$D_k$, and by \ref{C1} and the argument principle it has exactly one zero in
$D_k$, counted with multiplicity. Since $r_k<t_k$, Corollary~\ref{cor:zeta}
shows that this zero, call it $\rho_k$, is simple and lies on the critical
line.

\emph{The zeros are distinct and lie in $R$.} By \ref{C3} the discs are
pairwise disjoint: consecutive centres are further apart than the sum of the
radii, and the centres are collinear and ordered. Each $D_k$ lies in the open
upper half-plane, in the strip $-0.1\le\sigma\le1.1$, and below height
$2000$, hence in the interior of $R$. Because
$\zeta(\conj{s})=\conj{\zeta(s)}$, the conjugates $\conj{\rho_k}$ are simple
zeros lying in the reflected discs $\conj{D_k}$, which are in the open lower
half-plane and in the interior of $R$. The $2N=3034$ discs $D_k$,
$\conj{D_k}$ are pairwise disjoint, so $\zeta$ has at least $3034$ distinct
zeros in $R$.

\emph{There are no others.} The function $\zeta$ is meromorphic on a
neighbourhood of $R$ with exactly one pole inside, the simple pole at $s=1$.
By \ref{C2} and the argument principle, the number of zeros of $\zeta$ in $R$,
counted with multiplicity, is $3033+1=3034$. The $3034$ simple zeros found
above therefore exhaust them. All are simple and on the critical line.

\emph{The last sentence.} Every zero of $\zeta$ with non-zero imaginary part
satisfies $0\le\re s\le1$, so those with $0<\im s<2000$ lie in $R$. They are
the zeros in $R$ with positive imaginary part, namely $\rho_1,\dots,\rho_N$.
\qed

\begin{remark}
The proof never uses that $\zeta$ has no real zeros in $[-1,2]$. That fact is
a consequence: a real zero in $R$ would be a $3035$th.
\end{remark}

\begin{remark}
The count $1517$ is not an input. It is forced by \ref{C2}: if a zero had
been missed, or if a disc had held two, the totals $2N$ and $3034$ would not
agree and no conclusion would be drawn.
\end{remark}

\section{What must be trusted}\label{sec:trust}

The mathematical argument above is complete. The computed facts
\ref{C1}--\ref{C3} depend on the following.
\begin{enumerate}[leftmargin=2em]
\item \textbf{Ball arithmetic.} Enclosures are computed with Arb
  \cite{JohanssonArb}, through its Python binding \path{python-flint}
  $0.9.0$, at $128$ bits. In particular Arb's routine for $\zeta$ on a ball
  (see \cite{JohanssonHurwitz} for the underlying error bounds) must return a
  ball containing $\zeta(z)$ for every $z$ in the input ball, as must the
  routines for division, the complex exponential and the argument. The
  argument of this note does not depend on which algorithm Arb uses
  internally to evaluate $\zeta$, only on the enclosure property.
\item \textbf{The driver.} A script of about $250$ lines implements the
  procedure of Section~\ref{sec:winding}. Its logic is the four steps listed
  there.
\item \textbf{The geometry check.} Fact \ref{C3} is verified by a separate
  short script, named below, in exact rational arithmetic on the recorded
  doubles. It also checks that every recorded
  winding number is $1$ and that the rectangle's is $3033$.
\end{enumerate}
Neither Arb nor the script has been formally verified, and the certificate has
not been checked in a proof assistant. The result should be read as a
computer-assisted proof in the usual sense of interval arithmetic, not as a
formal one.

\paragraph{Reproduction.}
\sloppy
The script, its output and this note are in the repository
\url{https://github.com/judegomila/zeta-explorer}, under
\path{observatory/}: the script is
\path{scripts/experiments/zero-line-certificate.py} and the full record of
discs, windings and piece counts is
\path{dist/experiments/zero-line-certificate.json}; the geometry check is
\path{scripts/experiments/check-certificate-geometry.py}. The record carries the
SHA-256 digest of the script that produced it and the identifier
\path{zeta-zeros-below-2000-arb-v2}.

\section{Limits and remarks}\label{sec:limits}

\paragraph{Finite height.}
Theorem~\ref{thm:main} says nothing about zeros with $|\im s|\ge2000$. No
finite computation of this kind can prove the Riemann hypothesis.

\paragraph{Cost.}
The method is not a substitute for Turing's. It evaluates $\zeta$ on
two-dimensional boxes along contours that pass near zeros, where the image of
a box winds quickly and many bisections are needed. The number of boundary
pieces per disc in Table~\ref{tab:stats} already shows this. Methods based on
$Z(t)$ work on the line alone and scale to heights twelve orders of magnitude
beyond ours.

\paragraph{Close zeros.}
The radii shrink with the gaps between zeros. The smallest gap below height
$2000$ is about $0.0975$, giving a smallest radius of about $0.039$. At
greater heights, close pairs of zeros would force very small discs and
correspondingly fine subdivisions.

\paragraph{Other symmetric functions.}
Lemma~\ref{lem:trap} uses only the symmetry \eqref{eq:sym}. It therefore
applies to any holomorphic function with that symmetry, for example to the
functions obtained from $\xi$ by the de Bruijn--Newman heat flow, for which
the question of when all zeros lie on the line is the question of the sign of
the de Bruijn--Newman constant \cite{RodgersTao}. This note certifies nothing
about those functions.

\paragraph{Height $2000$.}
The height was chosen because it is the limit of the evaluator in the
accompanying viewer, so that every zero the viewer can display is covered by
a recorded disc.

\begin{thebibliography}{9}

\bibitem{Brent}
R.~P. Brent,
\emph{On the zeros of the Riemann zeta function in the critical strip},
Mathematics of Computation \textbf{33} (1979), 1361--1372.

\bibitem{JohanssonHurwitz}
F.~Johansson,
\emph{Rigorous high-precision computation of the Hurwitz zeta function and
its derivatives},
Numerical Algorithms \textbf{69} (2015), 253--270.

\bibitem{JohanssonArb}
F.~Johansson,
\emph{Arb: efficient arbitrary-precision midpoint-radius interval arithmetic},
IEEE Transactions on Computers \textbf{66} (2017), 1281--1292.

\bibitem{PlattTrudgian}
D.~Platt and T.~Trudgian,
\emph{The Riemann hypothesis is true up to $3\cdot10^{12}$},
Bulletin of the London Mathematical Society \textbf{53} (2021), 792--797.
arXiv:2004.09765.

\bibitem{RodgersTao}
B.~Rodgers and T.~Tao,
\emph{The de Bruijn--Newman constant is non-negative},
Forum of Mathematics, Pi \textbf{8} (2020), e6.
arXiv:1801.05914.

\bibitem{Turing}
A.~M. Turing,
\emph{Some calculations of the Riemann zeta-function},
Proceedings of the London Mathematical Society (3) \textbf{3} (1953), 99--117.

\end{thebibliography}

\end{document}
